\addtocontents{toc}{\protect\Needspace{3\baselineskip}}
\originalchapter{2004 年历史测验及闭合条件勘核}
B. Zwiebach, 2004 年 3 月 18 日, 限时 90 分钟, 允许个人两页笔记. 原两页扫描卷已逐页视觉读取公式, 不以有缺字的纯文本提取替代题面. 以下为四题完整重编与独立解答. 第四题最后的曲线族另需全局闭合条件, 本稿将其明确列出.

\src{8.251 Historical2004 Q1, (a)--(d), 各 5 分, 合计 20 分.}
\begin{exercise}
(a) 给 $\hbar$ 的 $M,L,T$ 量纲, 说明为何它在不同的时空维数下相同.
(b) 对 $S=-\frac12\int\dd^Dx\,\eta^{\mu\nu}\partial_\mu\phi\partial_\nu\phi$ 作 $\delta\phi$, 忽略边界项, 写成 $\delta S=\int\dd^Dx\,\delta\phi(\cdots)$, 求括号.
(c) 四维两静止 D1 膜均沿 $x^3$, 第一膜 $x^1=x^2=0$, 第二膜 $x^1=0,x^2=a>0$. 弦从第一膜（$\sigma=0$）到第二膜（$\sigma=\sigma_1$）. 绘图并列两端 $X^0,X^1,X^2,X^3$ 的八个条件.
(d) 用 $\mathcal L=-\frac{T_0}{c}\sqrt{(\dot X\cdot X')^2-\dot X^2X'^2}$ 与
$P^\sigma_\mu=\partial\mathcal L/\partial X'^\mu$, 求 $P^\sigma_\mu\dot X^\mu$.
\end{exercise}
\begin{solution}
(a) $[\hbar]=ML^2T^{-1}$, 等于作用量量纲, 因量子相位 $S/\hbar$ 必无量纲. 不同时空维数改变场的量纲和拉格朗日密度, 其积分仍为能量乘时间, 不改变作用量单位.

(b) 度规对称, 两个导数项合并后
$\delta S=-\int\partial^\mu\phi\,\partial_\mu\delta\phi\,\dd^Dx
=\int\delta\phi\,\partial^2\phi\,\dd^Dx$,
所以括号为 $\eta^{\mu\nu}\partial_\mu\partial_\nu\phi$.

(c) 每端的时间和 $X^3$ 均为自由, 即 $P_0^\sigma=P_3^\sigma=0$;
$X^1(t,0)=X^1(t,\sigma_1)=0$, $X^2(t,0)=0$, $X^2(t,\sigma_1)=a$ 均为 Dirichlet. 这列出四个自由条件与四个固定值.
\begin{center}
\begin{tikzpicture}[scale=.8]
 \draw[thick](0,-1)--(0,1.5)node[above]{第一 D1};
 \draw[thick](2,-1)--(2,1.5)node[above]{第二 D1};
 \draw[very thick,structurecolor](0,.3)..controls(.8,1)and(1.3,-.3)..(2,.5);
 \node[below]at(1,-1){沿 $x^2$ 相距 $a$, 均沿 $x^3$ 延伸};
\end{tikzpicture}
\end{center}

(d) 令 $\mathcal D=\sqrt{(\dot X\cdot X')^2-\dot X^2X'^2}$, 求偏导得
$P_\mu^\sigma=(T_0/c\mathcal D)[\dot X^2X'_\mu-(\dot X\cdot X')\dot X_\mu]$.
再与 $\dot X^\mu$ 缩并, 两项相消, 结果为零. 这是规范结构的恒等式, 无需假定 $\dot X\cdot X'=0$ 或先用运动方程.
\end{solution}

\src{8.251 Historical2004 Q2, (a)--(c), 20 分.}
\begin{exercise}
开弦能量 $E$, 在平面刚性旋转, 长度 $\ell=2E/(\pi T_0)$.
(a) 求 $\omega(E,T_0,c)$.
(b) 求从中点到一个自由端点的半根弦所携相对论总动量大小 $p$.
(c) 一个无质量粒子携带弦能量的一半, 其动量为 $q$, 求 $p/q$.
\end{exercise}
\begin{solution}
自由端点速度 $c$ 给 $\omega\ell/2=c$, 故 $\omega=\pi T_0c/E$.
一半直弦各点速度方向相同, 可将动量大小积分:
\[
 p=\int_0^{c/\omega}
 \frac{T_0\omega r\dd r}{c^2\sqrt{1-\omega^2r^2/c^2}}
 =\frac{T_0}{\omega}\int_0^1\frac{u\dd u}{\sqrt{1-u^2}}
 =\frac{E}{\pi c}.
\]
无质量粒子 $q=E/(2c)$, 因而 $p/q=2/\pi<1$.
两半弦总动量方向相反, 整弦在质心系总动量为零; (b) 问的是半根弦, 不能用整弦总动量替代.
\end{solution}

\src{8.251 Historical2004 Q3, (a)--(c), 20 分.}
\begin{exercise}
复平面先识别 $z\sim T_1z=z+1$, $z\sim T_2z=z+\tau_T$, $\tau_T=\ee^{\ii\pi/3}$.
基本平行四边形已有顶点 $0,1,\tau_T$.
(a) 给第四点、绘图、标边界识别并描述空间.
(b) 再加 $Rz=\ee^{2\pi\ii/3}z$. 除原点轨道外的两个固定点 $z_1,z_2$ 位于长对角线, $|z_1|<|z_2|$, 求之.
(c) $T_1\circ R$ 是否固定 $z_1$ 或 $z_2$? 区分平面代表与环面上的识别.
\end{exercise}
\begin{solution}
(a) 第四顶点为 $1+\tau_T$. 相对边通过 $T_1,T_2$ 平移配对, 得平坦环面 $\mathbb R^2/\Lambda$, $\Lambda=\mathbb Z+\mathbb Z\tau_T$.

(b) 因 $\tau_T^2=\tau_T-1$, $R=\tau_T-1$, 旋转保持该格, 所以在环面上良定义. 固定条件是 $(R-1)z\in\Lambda$, 不要求每个平面代表都满足 $Rz=z$. 长对角线上的解为
\[
 z_1=\frac{1+\tau_T}{3},\qquad
 z_2=\frac{2(1+\tau_T)}3,
 \quad Rz_1=z_1-1,\quad Rz_2=z_2-2.
\]
这两点与原点构成三种不同轨道. 也可用 $R-1$ 在格基中的整数矩阵行列式绝对值 3 核对固定点总数.

\begin{center}
\begin{tikzpicture}[scale=2.6,>=Stealth]
 \coordinate(O)at(0,0);\coordinate(A)at(1,0);
 \coordinate(B)at(.5,.8660254);\coordinate(C)at(1.5,.8660254);
 \draw[thick](O)--(A)--(C)--(B)--cycle;
 \draw[->](.25,0)--(.65,0);\draw[->](.75,.866)--(1.15,.866);
 \draw[->](.15,.2598)--(.35,.6062);\draw[->](1.15,.2598)--(1.35,.6062);
 \draw[dashed](O)--(C);
 \fill(.5,.288675)circle(.7pt)node[below]{$z_1$};
 \fill(1,.577350)circle(.7pt)node[above]{$z_2$};
 \node[below]at(O){$0$};\node[below]at(A){$1$};
 \node[above]at(B){$\tau_T$};\node[above]at(C){$1+\tau_T$};
\end{tikzpicture}
\end{center}

(c) 作为平面点, $(T_1R)z_1=z_1$, $(T_1R)z_2=z_2-1\ne z_2$, 所以只有 $z_1$ 被这个特定平面变换严格固定.
但在环面上平移 $T_1$ 已是识别, 两者仍都是固定轨道. 此区别避免把基本域边界代表与商空间的点混为一谈.
\end{solution}

\src{8.251 Historical2004 Q4, (a)--(g), 40 分.}
\begin{exercise}
静态规范下闭弦满足
$\boldsymbol X_{\sigma\sigma}-c^{-2}\boldsymbol X_{tt}=0$,
$|\boldsymbol X_\sigma\pm c^{-1}\boldsymbol X_t|^2=1$, $\sigma\sim\sigma+L$.
(a) 初始 $\boldsymbol X_t(0,\sigma)=0$, 用一个向量函数 $\boldsymbol F(u)$ 写通解.
(b) 给参数约束对 $\boldsymbol F$ 的条件.
(c) 闭弦周期性对 $\boldsymbol X$ 和 $\boldsymbol F$ 有何要求?
(d) 若 $\boldsymbol F'=(-\sin(2\pi u/L),\cos(2\pi u/L))$, 描述 $t=0$ 构型.
以下再取
\[
 \boldsymbol F'(u)=
 \left(-\sin[\gamma\sin(2\pi u/L)],
 \cos[\gamma\sin(2\pi u/L)]\right),\qquad\gamma>\pi/2.
\]
(e) 化简 $|\boldsymbol X_t|/c$.
(f) 给突然首次出现尖点的最小时间 $t_*>0$ 的方程.
(g) 若 $3\pi/2<\gamma<2\pi$, 同时最多有几个尖点?
\end{exercise}
\begin{solution}
(a) 零初速度使反向两行波相等, 故
$\boldsymbol X(t,\sigma)=\frac12[\boldsymbol F(\sigma+ct)+\boldsymbol F(\sigma-ct)]$.
初始构型就是 $\boldsymbol F(\sigma)$, 可吸收常向量平移.
(b) 约束等价于 $|\boldsymbol F'(u)|=1$.
(c) $\boldsymbol X(t,\sigma+L)=\boldsymbol X(t,\sigma)$, 特别在 $t=0$ 必有 $\boldsymbol F(u+L)=\boldsymbol F(u)$, 即不仅要求 $\boldsymbol F'$ 周期, 还要求 $\int_0^L\boldsymbol F'\dd u=0$.
(d) 积分得 $\boldsymbol F=(L/2\pi)(\cos(2\pi u/L),\sin(2\pi u/L))+\text{常向量}$, 所以初始是半径 $L/(2\pi)$ 的静止圆.

(e) 令 $S=2\pi\sigma/L,T=2\pi ct/L$. 速度是两个单位切向量之差的一半, 两向量夹角的半差为
$\gamma[\sin(S+T)-\sin(S-T)]/2=\gamma\cos S\sin T$. 因而
\[
 \frac{|\boldsymbol X_t|}{c}
 =|\sin(\gamma\cos S\sin T)|.
\]
(f) 尖点条件为此速度达到 1（同处空间切向导数为零）. 最初 $\sin T$ 从零增加, 其最大角幅首次达到 $\pi/2$ 时出现:
\[
 \gamma\sin\frac{2\pi ct_*}{L}=\frac\pi2,
 \quad t_*=\frac{L}{2\pi c}\arcsin\frac{\pi}{2\gamma}.
\]
最初发生在 $\cos S=\pm1$, 即 $\sigma=0,L/2$, 零与 $L$ 是同一点而不重复计数.
(g) 在给定 $\gamma$ 范围, 若取 $|\sin T|=1$, 允许角值为 $\pm\pi/2,\pm3\pi/2$. 每个值都严格在振幅区间内, 各有两个不同的 $S$ 解, 最多共 8 个. 不会出现 $\pm5\pi/2$, 因 $\gamma<2\pi$.

必须另核对 (c). 对最后的切向量族, 第一分量整圈平均为零, 第二分量平均为
\[
 J_0(\gamma)=\frac1{2\pi}\int_0^{2\pi}\cos(\gamma\sin s)\dd s.
\]
这正是零阶 Bessel 函数的一种积分定义. 所以
$\boldsymbol F(u+L)-\boldsymbol F(u)=L(0,J_0(\gamma))$.
原题只给 $\gamma>\pi/2$ 或 (g) 的区间, 没有要求 $J_0(\gamma)=0$, 因而任意这样的 $\gamma$ 一般不是闭合弦. (e)--(g) 是该局部波动解的正确尖点计算; 要同时满足原题“闭弦”的全局条件, 还须选择 $J_0(\gamma)=0$.
例如 (g) 区间内可取数值根 $\gamma\simeq5.52007811$, 此时构型闭合且上述最大 8 个尖点确实兼容. 本稿保留并指出该缺项, 不暗中将所有 $\gamma$ 写成已闭合的曲线.
\end{solution}

